% Session 6 -- The Cover Sheet and the Records Retained: Procedures TPS-P-06
% and TPS-P-07, Checklist TPS-C-2.
% Law 1 sources: spec/tps_ritual.json R5 (Y continues to the report; N ends
% with return code 5 and no report file; other keys ignored) and file_rules
% (TPSREPT.BAK deleted, TPSREPT.TXT renamed TPSREPT.BAK, TPSREPT.TXT
% created and written, Ctrl-Z padded); os/milton/tpsrep_program.asm .r5
% (the key folded to upper case, echoed, no ENTER), report (m_done1, and
% m_done2 only when there was a previous report), totals (the corrective
% action line only when a report lacks its cover sheet).  Screens tpsdone,
% tpsbak: the journey route; test-flair-tps-journey and test-tps-report
% prove TPSREPT.TXT byte for byte against real MS-DOS 6.22.
% The R5 prompt names FORM TPS-1 (REV. 10/89) since lane tpsr2 (77835faa,
% bead initech-5igf).  Screens r5cover (the run waiting at the prompt: the
% answer Y scrolls it away with the report) and r5refuse: the journey route.
% GATES: test-tps-report test-flair-tps-journey test-pubs-screens test-pubs-tps
\Gates{test-tps-report, test-flair-tps-journey, test-pubs-screens, test-pubs-tps}
\chapter{The Cover Sheet and the Output}

\SideHead{Introduction}When the memo has been acknowledged, the program asks
whether the new cover sheet is attached.  Only when the operator answers
\type{Y} does it write the report.  In this session you attest the cover
sheet, verify the report, and retain the records of the run.

\SideHead{Objectives}Upon completion of this session you will be able to:
\begin{enumerate}
\item Attach the new cover sheet, \form{1} (Rev.\ 10/89), and attest it.
\item State what happens when the answer is \type{N}.
\item Verify that the report has been written and the previous report kept.
\item Complete Checklist TPS-C-2 and \form{5}.
\end{enumerate}

\section{The Cover Sheet Attestation}
\index{cover sheet attestation}\index{Form TPS-1}\index{TPS-005@\texttt{TPS-005}}
The program asks:
\TpsMsg{R5.prompt}
The answer is one key, \type{Y} or \type{N}, in capital or small letters;
\key{Enter} is not pressed.  Any other key is ignored.  The program displays
the answer as a capital letter.

\SideHead{Y}The program writes the report (\proc{07}).

\SideHead{N}The program displays:
\TpsMsg{TPS-005}
where \ph{n} is the memo number, and ends with return code 5.  No report is
written, and \rpt\ and \file{TPSREPT.BAK} are left as they were.

\begin{Procedure}{TPS-P-06}{To Attest the Cover Sheet}
\ProcField{Purpose}{To attach the new cover sheet to the report folder and
  attest it to the program.}
\ProcField{Frequency}{Each run, directly after \proc{05}.}
\ProcField{Performed by}{The operator.}
\ProcField{Prerequisites}{The cover sheet prompt is displayed.}
\ProcField{Materials}{A new cover sheet, \form{1} (Rev.\ 10/89); the report
  folder for this run.}
\begin{ProcSteps}
\item Check that the cover sheet is \form{1} with the revision date 10/89.
  An earlier revision is not the new cover sheet.  If you do not have one,
  answer \type{N} and see your supervisor.
\item Enter on \form{1} the report number from \form{2}, column 2, and the
  date.  Attach it to the front of the report folder.
\item At the cover sheet prompt, press \type{Y}.
  \OpsShot[fromtext={KEY THE NUMBER}]{r5cover}
\end{ProcSteps}
\ProcField{Verification}{The report begins to be displayed.}
\ProcField{Records retained}{\form{1}, in the report folder.}
\ProcNote{An answer of \type{N} cannot be taken back.  To write the report,
  run the program again from \proc{02}.}
\end{Procedure}

\OpsShot[fromtext={DID YOU REMEMBER},
  caption={The answer N: the report is not written.}]{r5refuse}

\section{The Report}
\index{TPSREPT.TXT}\index{TPSREPT.BAK}
When the answer is \type{Y}, the program keeps one previous report and
writes the new one:
\begin{enumerate}
\item It deletes \file{TPSREPT.BAK}, if there is one.
\item It renames \rpt, if there is one, to \file{TPSREPT.BAK}.
\item It creates \rpt\ and writes the report to it, displaying each line on
  the screen as it is written: the cover sheet, the detail pages and the
  totals (Session~1).
\end{enumerate}
It then displays \texttt{TPS REPORT WRITTEN TO TPSREPT.TXT.}  When there was
a previous report, it also displays \texttt{PREVIOUS REPORT RETAINED AS
TPSREPT.BAK.}  It ends with return code 0.

\begin{Procedure}{TPS-P-07}{To Verify and Retain the Output}
\ProcField{Purpose}{To verify that the report has been written and the
  previous report kept, and to retain the records of the run.}
\ProcField{Frequency}{Each run that ends with return code 0.}
\ProcField{Performed by}{The operator.}
\ProcField{Prerequisites}{\proc{06} is complete with the answer \type{Y}.}
\ProcField{Materials}{\form{5}; Checklist TPS-C-2.}
\begin{ProcSteps}
\item Wait until the prompt \texttt{A:\bs APPS>} is displayed.  Do not press
  any key while the report is being written.
\item Read the last lines of the screen.
  \OpsShot[fromtext={REPORTS LISTED}, marks={1=WITHOUT COVER;2=CORRECTIVE;3=TPS REPORT WRITTEN;4=TPS REPORT GENERATOR TERMINATED},
    label=fig:done, caption={The end of the report and of the run.}]{tpsdone}
\item Enter the figure \texttt{WITHOUT COVER SHEET} on \form{5}, line 7.  If
  the line \texttt{CORRECTIVE ACTION REQUIRED} is displayed, enter it too,
  and give the report folder to your supervisor with the memo it names.
\item Type \type{DIR TPSREPT.* /W} and press \key{Enter}.
  \OpsShot[fromtext={A:\Bsl APPS>DIR TPSREPT}]{tpsbak}
\item Complete Checklist TPS-C-2 and sign \form{5}.
\end{ProcSteps}
\ProcField{Verification}{\texttt{RETURN CODE 0} is displayed; the listing
  names \rpt\ and, after the first run, \file{TPSREPT.BAK}.}
\ProcField{Records retained}{\rpt\ (this report) and \file{TPSREPT.BAK} (the
  previous report) in \file{A:\bs APPS}; \form{1} through \form{5};
  Checklists TPS-C-1 and TPS-C-2.}
\ProcNote{The program keeps one previous report only.  The report that was
  in \file{TPSREPT.BAK} is deleted by the next run (Checklist TPS-C-1).}
\end{Procedure}

\begin{CalloutLegend}
\LegendRow{1}{Without cover sheet}{The number of TPS reports on file filed
  without a cover sheet.  Entered on \form{5}, line 7.}
\LegendRow{2}{Corrective action}{Displayed only when that number is not
  zero.  Names the memo.}
\LegendRow{3}{Report written}{\rpt\ has been written.}
\LegendRow{4}{End line}{Return code 0: the run is complete.}
\end{CalloutLegend}

When the run is the second or a later one, the end of the screen shows the
previous report retained:

\OpsShot[fromtext={TPS REPORT WRITTEN}, totext={TPS REPORT GENERATOR TERMINATED},
  caption={A second run: the previous report is retained.}]{tpsbak}

\section{After the Run}
\index{post-run checklist}\index{operator sign-off sheet}\index{Form TPS-5}
\begin{Checklist}{TPS-C-2}{Post-Run}
\CheckItem{\texttt{RETURN CODE 0} was displayed.}
\CheckItem{\rpt\ is listed in \file{A:\bs APPS}.}
\CheckItem{\form{1} (Rev.\ 10/89) is attached to the front of the report
  folder.}
\CheckItem{\form{2}, \form{3} and \form{4} carry this run's line.}
\CheckItem{The figure \texttt{WITHOUT COVER SHEET} is entered on \form{5},
  line 7, and any corrective action has been given to your supervisor.}
\CheckItem{\form{5} is signed and, with Checklists TPS-C-1 and TPS-C-2,
  filed in the report folder.}
\end{Checklist}

\section{Lab Project 5. The Cover Sheet and the Report}
Continue from Lab Project 4: the program displays the cover sheet prompt.
\begin{enumerate}
\item Carry out \proc{06} with a copy of \form{1} from Appendix F.
\item Carry out \proc{07}.  Compare the end of your screen with
  Figure~\ref{fig:done}.
\item Run the program a second time, from \proc{02}, with report number
  1015.  After \proc{07}, step 4, the listing names \file{TPSREPT.BAK}.
\end{enumerate}

\section{Study Questions}
\begin{enumerate}
\item What is the revision date of the new cover sheet?
\item What happens when the answer to the cover sheet prompt is \type{N}?
\item Is \key{Enter} pressed after the answer?
\item Which file holds the report written before this one?
\item When is the line \texttt{CORRECTIVE ACTION REQUIRED} displayed, and
  what do you do?
\end{enumerate}
